%%%%%%%%%%%%%%%%%%%%%%%%%%%%%%%%%%%%%%%
% Cover Letter - Coleman Group (методолог по управлению бизнес-процессами)
%%%%%%%%%%%%%%%%%%%%%%%%%%%%%%%%%%%%%%%

\documentclass[]{cover}
\usepackage{fancyhdr}

\pagestyle{fancy}
\fancyhf{}

\rfoot{Page \thepage \hspace{0pt}}
\thispagestyle{empty}
\renewcommand{\headrulewidth}{0pt}

\begin{document}

%%%%%%%%%%%%%%%%%%%%%%%%%%%%%%%%%%%%%%
%     TITLE NAME
%%%%%%%%%%%%%%%%%%%%%%%%%%%%%%%%%%%%%%
\namesection{}{\Huge{Алексей Орлов}}{  \href{mailto:lex.orlov78@yandex.ru}{lex.orlov78@yandex.ru} | +7~(901)~797-54-88 |  \urlstyle{same}\href{https://www.linkedin.com/in/alex-orlov-v/}{LinkedIn}
}

%%%%%%%%%%%%%%%%%%%%%%%%%%%%%%%%%%%%%%
%     MAIN COVER LETTER CONTENT
%%%%%%%%%%%%%%%%%%%%%%%%%%%%%%%%%%%%%%

\currentdate{\rutoday}
\lettercontent{Здравствуйте, команда Coleman Group,}

\lettercontent{Меня заинтересовала позиция методолога по управлению бизнес-процессами. Сразу хочу быть честным: формальной нотации BPMN 2.0 и стандарта BPM CBOK 4.0, которые вакансия называет обязательными, у меня нет. Есть около 6 лет практики построения методологии, целевых моделей процессов и регламентирующей документации на уровне нескольких направлений одновременно, сначала в Банке ВТБ, затем в ООО «ПочтаТех», и ниже я показываю, где этот опыт закрывает суть задачи даже без формальной нотации.}

\lettercontent{Что я принесу на позицию методолога:}

{\raggedright\fontspec{DejaVu Sans}\fontsize{11pt}{13pt}\selectfont
\begin{itemize}
    \item \textbf{Целевая модель и регламентация:} описал целевую модель процессов (To-Be) и сравнил её с текущей (As-Is) для нескольких ИТ-направлений ПочтаТех, выявляя разрывы и дублирование действий
    \item \textbf{Governance на уровне топ-менеджмента:} согласовывал целевую модель процесса с CTO и руководителями направлений в рамках квартального планирования
    \item \textbf{Обучение и менторство:} запустил 2 внутренние образовательные программы (Школу Scrum-мастеров и Школу Agile-коучей) и обучил 50+ сотрудников в Банке ВТБ; провёл ассесмент 25+ сотрудников (PM/Scrum Master/Delivery Manager) с разработкой треков развития в ООО «ПочтаТех»
    \item \textbf{Архитектура на стыке направлений:} выстроил единый процесс управления изменениями, затрагивающий несколько бизнес-направлений и ИТ-команд одновременно
\end{itemize}\par}
\vspace{6pt}

\lettercontent{Формальную нотацию BPMN и CBOK беру на себя как задачу первых недель: методологическая база и опыт кросс-функциональной трансформации у меня уже есть, конкретный инструмент поверх неё осваивается предметно. Буду рад обсудить детали.}

\begin{flushright}
\closing{С уважением,}

\signature{Алексей Орлов}
\end{flushright}
\end{document}
